% problem_id: S014-lurie-ultracategories-theorem-2-1-1
% source_id: S014 (Jacob Lurie, Ultracategories)
% source_file: lurie-conceptual.pdf
% statement_locator: printed p. 22
% proof_locator: printed p. [22]
% representation: PDF; transcribed from rendered page images (never from the text layer)
% status: complete (statement + author proof verbatim + reason + locators)
%
%% problem_id: L1
%% source_locator: J. Lurie, "Ultracategories", Theorem 2.1.1 (Los Ultraproduct Theorem),
%%   section 2.1 "The Los Ultraproduct Theorem", printed page 22 = PDF page 22.
%%   Statement and proof are on the SAME page (printed/PDF page 22); the proof does NOT cross
%%   to page 23 (page 23 begins with "Corollary 2.1.4").
%% source_pdf: /disks/sata1/yupeng/human-proof-corpus/source-cache/registry-sources/lurie/lurie-conceptual.pdf
%% method: pages located with `pdftotext -layout | grep -n` ONLY; all text below transcribed
%%   by eye from rendered page images (pdftoppm -r 150 -png, proof paragraph and the composition
%%   diagram re-rendered at -r 400 for glyph-level checking). The text layer of this work is
%%   unreliable (it renders "∈" as "2" and "α" as "↵") and was never used as source text.
%% conventions: no rewording, no simplification, no completion. Original typography kept.
%%   Suspect spots in the printed original are copied as printed and flagged in a comment,
%%   never silently corrected.

%% --- page 22, section 2.1 lead-in (context that fixes the notation used by the statement;
%%     this paragraph is NOT part of Theorem 2.1.1) ---

Note that, for any category $\mathcal{C}$, the functor category $\mathrm{Fun}(\mathcal{C},\mathrm{Set})$
admits small limits and colimits. In particular, for any collection of functors
$\{M_s : \mathcal{C} \to \mathrm{Set}\}_{s \in S}$ and any ultrafilter $\mu$ on $S$, we can form
the categorical ultraproduct $\int_S M_s d\mu$ (in the category $\mathrm{Fun}(\mathcal{C},\mathrm{Set})$),
which is described concretely by the formula
\[
\Big(\int_S M_s d\mu\Big)(C) \;=\; \int_S (M_s(C)) d\mu \;=\; \varinjlim_{\mu(S_0)=1}\ \prod_{s \in S_0} M_s(C).
\]

%% --- page 22, statement ---

\begin{theorem}[Los Ultraproduct Theorem]
%% printed: "Theorem 2.1.1 (Los Ultraproduct Theorem)."
Let $\mathcal{C}$ be a pretopos and let $\{M_s\}_{s \in S}$ be a collection of models
of $\mathcal{C}$. For every ultrafilter $\mu$ on $S$, the ultraproduct $\int_S M_s d\mu$
(formed in the category $\mathrm{Fun}(\mathcal{C},\mathrm{Set})$) is also a
model of $\mathcal{C}$.
\end{theorem}

%% --- page 22, proof ---

\noindent\textit{Proof of Theorem 2.1.1.} Let $\{M_s\}_{s\in S}$ be a collection of models of
$\mathcal{C}$ and let $\mu$ be an ultrafilter on the index
set $S$. Then the ultraproduct $\int_S M_s d\mu$ (formed in the category
$\mathrm{Fun}(\mathcal{C},\mathrm{Set})$) can be identified with the
composition
\[
\mathcal{C} \xrightarrow{\ \{M_s\}_{s\in S}\ } \mathrm{Set}^{S}
            \xrightarrow{\ \int_S(\bullet)d\mu\ } \mathrm{Set}.
\]
The first map is a pretoos functor by virtue of our assumption that each $M_s$ is a model of
$\mathcal{C}$, and the second
map is a pretopos functor by virtue of Proposition 2.1.3. It follows that the composite map is
also a pretopos
functor. \hfill$\square$

%% ---------------------------------------------------------------------------
%% TRANSCRIBER'S NOTES (nothing above has been altered)
%%
%% [ORIGINAL TYPO] In the proof, the first occurrence reads "pretoos functor" — sic. The word
%%   is printed exactly as "pretoos" (missing the "p" of "pretopos"); the very next sentence
%%   spells the same word correctly as "pretopos". Verified at 400 dpi. Left verbatim above;
%%   NOT corrected. (This typo is also present in the PDF text layer, which reads
%%   "The first map is a pretoos functor".)
%%
%% [NO OTHER ANOMALIES] No repeated words, missing words, or illegible glyphs were found in the
%%   statement or the proof. No %% [ILLEGIBLE: ...] markers were needed.
%%
%% [HYPERREF ARTIFACTS] In the rendered original, "2.1.1", "Proposition 2.1.3", "1.2.2", "2.3.1",
%%   "2.4.2", "1.3.7", "§2.3", "§2.4" appear inside red hyperlink boxes; these boxes are
%%   hyperref decorations, not mathematical content, and are omitted above.
%%
%% [DIAGRAM] The composition diagram is printed with the arrow labels ABOVE the arrows:
%%   "{M_s}_{s∈S}" over the first arrow and "∫_S(●)dμ" over the second. Reproduced faithfully.
%%
%% [STRUCTURE OF THE PROOF — for the difficulty judgement, not part of the transcription]
%%   The whole proof is this one paragraph: it (a) identifies ∫_S M_s dµ with a composite of two
%%   functors, (b) cites the hypothesis that every M_s is a model of C for the first functor,
%%   (c) cites Proposition 2.1.3 for the second functor, (d) invokes closure of pretopos functors
%%   under composition. The substantive content is entirely delegated to Proposition 2.1.3,
%%   proved separately just above on the same page.
%% ---------------------------------------------------------------------------
